\documentclass[10pt]{article}
\usepackage[margin=0.8in]{geometry}
\usepackage{amsmath,amssymb,amsthm,enumitem}
\usepackage[hidelinks]{hyperref}
\usepackage{xurl}
\newcommand{\Tr}{\operatorname{Tr}}
\newcommand{\tr}{\operatorname{tr}}
\newcommand{\id}{\mathbf 1}
\newcommand{\dd}{\,\mathrm d}
\begin{document}
\begin{center}\textbf{The graded planar cut of the four-packet sine sea}\end{center}

\paragraph{Question and separation of sources.}
This derivation concerns the native site algebra on a flat slice,
with one planar interface. It tests the signed combination
$S_{\rm gr}=\sum_s\eta_s S_s$; it does not identify that combination
with the von Neumann entropy of a physical density matrix or with a
thermodynamic entropy. The four-packet Hamiltonian, its positive
excitation sector and its vacuum coefficient are fixed as in Sources~1--3.
No derivation, route or numerical output of the icloud Strand S is
used. No Pauli--Villars induced-gravity entropy formula is an input.

\paragraph{Analytic result.}
For the positive rectangle below, the native flat planar cut obeys
\begin{equation}
 \frac{S_{\rm gr}}{\mathcal A}
 \longrightarrow \frac{J}{3\pi},\qquad
 \frac{S_{\rm gr}/\mathcal A}{(\pi/2)k_a(1)}
 \longrightarrow \frac29 .
 \label{eq:s-claim}
\end{equation}
The exact half-chain Toeplitz/CTM reduction and a grouped continuum
cone calculation give the same coefficient. The proposed equality
to $\mathcal A/(4G)$ fails for all sufficiently small $a$ at each
fixed positive rectangle. This does not exclude an incidental
equality at another finite spacing.
% STATUS [X]: exact Hardy/Toeplitz/CTM reduction and shifted lattice measure;
% independently normalized grouped Dirac-cone route. No [R] or [N] claim.

\section*{1. Exact cut, sea and multiplicities}
Use $\hbar=c=k_B=1$ and natural logarithms. Set $N=\psi=1$
and choose the single three-dimensional sine operator
\[
 D_a(p)=\sum_{i=1}^3\sigma_i\frac{\sin(ap_i)}a,
 \qquad A_s=m_s+iD_a,
\]
and fill the negative band of
\begin{equation}
 h_s(p)=m_s\beta+\sum_i\alpha_i\frac{\sin(ap_i)}a,
 \quad \beta=\tau_1\otimes\id_2,\quad
 \alpha_i=\tau_2\otimes\sigma_i.
 \label{eq:s-flat-h}
\end{equation}
There is one four-component lattice fermion per packet, with eight
spatial valleys, not eight independent lattice fields. Time is
continuous and supplies no temporal doubling. The squared masses are
\[
 (x_s)=(x_0,x_0+h_1,x_0+h_2,x_0+h_1+h_2),\qquad
 x_0,h_1,h_2>0,
 \quad\eta=(1,-1,-1,1),\quad m_s=\sqrt{x_s}.
\]
Thus $\sum_s\eta_s=\sum_s\eta_sx_s=0$ and
\begin{equation}
 J=\sum_s\eta_sx_s\log(x_s/\mu_0^2)
   =\int_0^{h_1}\!\dd u\int_0^{h_2}\!\dd v\,
       \frac1{x_0+u+v}>0 .
 \label{eq:s-J}
\end{equation}
The arbitrary logarithm scale cancels. In particular, the plus sign
in this definition agrees with the positive rectangle integral.

Use the CAR algebra generated by all native fermions at
$n_1\ge0$. The interface lies between lattice planes $n_1=-1$ and
$n_1=0$. It has physical area $\mathcal A$ at the unit ruler.
Tangential directions are periodic before their thermodynamic limit.
The normal ends are sent to infinity with $m_s L_1\to\infty$;
there is one interface, not a two-face slab. Take these large-volume
limits at fixed $a>0$ and positive masses before $a\to0$.
No curvature, corner or finite-radius contribution is included in the
planar area density.

The one-packet negative-band projector and restricted correlation are
\[
 P_s(p)=\frac12\left(\id_4-
       \frac{h_s(p)}{\sqrt{x_s+\sum_i\sin^2(ap_i)/a^2}}\right),
 \qquad C_s=\chi_{n_1\ge0}P_s\chi_{n_1\ge0}.
\]
For $f(c)=-c\log c-(1-c)\log(1-c)$,
$S_s=\Tr f(C_s)$, understood per unit interface area in the stated
limit. Reversing a whole packet replaces $h_s$ by
$\beta h_s\beta$ at this flat background. That site-local unitary
preserves every cut eigenvalue. Hence the signs in $S_{\rm gr}$ are
an additional linear grading of positive entropies, not signs
generated by packet reversal itself.

The protected sea of the completed light packet is the same free
sea, since the quartic excitation generator annihilates that state.
The empty number-conserving partner-boson vacuum is a product in
its native site oscillator algebra; it supplies no additional cut
entropy here. Its zero-point trace in $W_0$ is not a fermion entropy
multiplicity.
% SOURCES: 1--3; the algebra and the boundary convention above define this task.

\section*{2. Exact dimensional reduction before grading}
Fourier transformation in the two tangential directions acts
separately within the region and its complement. At fixed
$\theta_\perp=(ap_2,ap_3)$, the three anticommuting mass matrices
$m\beta+\alpha_2\sin\theta_2/a+\alpha_3\sin\theta_3/a$
have norm
\[
 M_\perp=\sqrt{m^2+a^{-2}
                    (\sin^2\theta_2+\sin^2\theta_3)}.
\]
A constant unitary rotation at that tangential momentum gives two
identical two-component chains with symbol
$a^{-1}(\sigma_3\sin\theta_1+\sigma_1 aM_\perp)$.
It is independent of $n_1$ and preserves the cut. Let $s_1(t)$ be
the one-interface entropy of one such two-component chain with
dimensionless mass $t>0$. Then the exact lattice expression is
\begin{equation}
 \frac{S_s(a)}{\mathcal A}
 =\frac2{a^2}\int_{[-\pi,\pi]^2}
       \frac{\dd^2\theta}{(2\pi)^2}\,
       s_1\!\left(\sqrt{\sin^2\theta_2+\sin^2\theta_3+a^2x_s}\right).
 \label{eq:s-exact}
\end{equation}
The factor two is the four-component spin-block multiplicity.
The two normal sine zeros remain inside $s_1$; the four tangential
zeros remain inside this integral. They together give eight
four-component continuum Dirac valleys.

Group the four terms in \eqref{eq:s-exact} before any expansion.
Writing its right side as $F_a(x_s)$ gives the exact finite-mass
rectangle identity
\begin{equation}
 \frac{S_{\rm gr}(a)}{\mathcal A}
 =\int_0^{h_1}\!\dd u\int_0^{h_2}\!\dd v\,
       F_a''(x_0+u+v).
 \label{eq:s-groupfirst}
\end{equation}
This equation fixes the order of grouping and limiting operations.
It does not presuppose that the finite-$a$ entropy is a heat trace.

\section*{3. Infrared coefficient and the lattice remainder}
One complex two-component Dirac field in one space dimension has
the one-interface logarithm $-(1/6)\log(M\epsilon)$. The
normal sine chain has two Dirac zeros. The needed asymptotic is
only
\begin{equation}
 s_1(\sqrt z)=-\frac16\log z+c_0+o(1),\qquad z\downarrow0.
 \label{eq:s-chainlog}
\end{equation}
No derivative bound for the critical entropy is needed below.
The logarithmic coefficient is obtained from a replica cone: for a
complex two-component Euclidean Dirac field,
$\log Z=-\tfrac12\int \dd t\,t^{-1}e^{-M^2t}
\Tr e^{-t(-\nabla^2+R/4)}$. The linear-curvature heat coefficient is
\[
 \tr_2(R/6-R/4)=-R/6,
 \qquad \int_{C_n}R=4\pi(1-n).
\]
Consequently
$\log Z_n-n\log Z_1=-(n-1)\int\dd t\,t^{-1}e^{-M^2t}/12$
to first order in $n-1$. Applying $-\partial_n|_{n=1}$ yields
$\tfrac1{12}\int_{\epsilon^2}^{\infty}\dd t\,t^{-1}e^{-M^2t}
=-(1/6)\log(M\epsilon)+O(1)$ for one Dirac zero.
This cone is the gluing of the flat density-matrix path integral;
it is not an assumption about the inhomogeneous polar law's
coupling to a physical conical metric.

\paragraph{Exact half-chain map into the CTM theorem's domain.}
A constant orbital rotation puts the flattened native chain in
off-diagonal form with phase
\[
 \phi_t(z)=\frac{t+(z-z^{-1})/2}{\sqrt{t^2+\sin^2\theta}},
 \qquad z=e^{i\theta}.
\]
Let $T_\phi=\Pi_+M_\phi\Pi_+$ be compression to the Hardy
space $H^2$, identified with the sites $n_1\ge0$. Its restricted
correlation is, up to a constant orbital rotation,
\begin{equation}
 C_t=\frac12
 \begin{pmatrix}\id&-T_{\phi_t}\\-T_{\phi_t}^{\dagger}&\id\end{pmatrix}.
 \label{eq:s-toeplitz-cut}
\end{equation}
For $r=\sqrt{1+t^2}-t\in(0,1)$ the phase satisfies
\[
 \phi_t(z)^2=-\frac{(z-r)(z+r^{-1})}{(z+r)(z-r^{-1})}.
\]
Choose any $b\in(r,1)$ and set
\[
 B_b(z)=\frac{z+b}{1+bz},\quad
 A=\frac{b-r}{1-br},\quad B=\frac{b+r}{1+br},\qquad
 0<A<B<1.
\]
The weighted composition
\[
 (W_bf)(z)=\frac{\sqrt{1-b^2}}{1+bz}f(B_b(z))
\]
is unitary on the circle: its squared weight is the angular
Jacobian of $B_b$. It maps $H^2$ onto $H^2$, since both the map
and its inverse are disk automorphisms with analytic nonvanishing
weights. It therefore also preserves the orthogonal cut complement.
Multiplication obeys
$W_bM_\phi W_b^{\dagger}=M_{\phi\circ B_b}$; using its inverse
shows that $T_{\phi_t}$ and $T_{\phi_t\circ B_b^{-1}}$ are
unitarily equivalent. This unitary need not be local along the
chain; preservation of the two cut subspaces is what is required.

The transformed phase has squared-phase zeros $B,A^{-1}$ and
poles $A,B^{-1}$. Define
\[
 h=\frac{1+AB}{A+B}>1,\qquad
 \gamma=\frac{B-A}{A+B}\in(0,1),\qquad
 G(z)=h-\cos\theta+i\gamma\sin\theta.
\]
The ratio $G/\overline G$ has exactly these zeros and poles.
Both normalized phases equal one at $z=1$, fixing the square-root
sign. Hence the transformed phase is $G/|G|$, equal up to a constant
minus sign to the XY phase in the primary source. The minus sign
does not change singular values. This puts the native compression
\emph{inside} the printed $h>1$, $0<\gamma<1$ CTM domain;
no large-anisotropy continuation is used.

If $\nu_j$ are the nontrivial singular values of this compression,
\eqref{eq:s-toeplitz-cut} has paired eigenvalues
$(1\pm\nu_j)/2$. Its entropy is twice the XY Majorana-covariance
entropy, whose sum contains one binary entropy per $\nu_j$.
For contiguous halves of an open chain, the Jordan--Wigner
occupation basis is the same site product basis; its spin and
regional CAR reduced spectra agree.
Peschel's one-interface spectrum is
$\varepsilon_j=(2j+1)\varepsilon$ for $h>1$ (Source~4).
Its modulus is
\[
 \frac{\gamma}{\sqrt{h^2+\gamma^2-1}}
 =\frac{B-A}{1-AB}=\frac{2r}{1+r^2}
 =\frac1{\sqrt{1+t^2}}.
\]
Thus the exact native half-chain entropy is
\begin{equation}
 k=(1+t^2)^{-1/2},\quad k'=t(1+t^2)^{-1/2},\quad
 \varepsilon=\pi\frac{K(k')}{K(k)},\quad
 s_1(t)=2\sum_{j=0}^{\infty}F((2j+1)\varepsilon),
 \label{eq:s-ctm}
\end{equation}
where $F(v)=\log(1+e^{-v})+v/(1+e^v)$ and $K$ is the complete
elliptic integral. This is a half-chain spectrum: no second interface
is included. Its finite critical constant follows directly:
\[
 \int_0^\infty F(v)\dd v=\frac{\pi^2}{6},\qquad
 2\sum_{j\ge0}F((2j+1)\varepsilon)
 =\frac{\pi^2}{6\varepsilon}+O(\varepsilon).
\]
The midpoint bound uses the integrable second derivative of $F$.
The elliptic asymptotics $K(k)=\log(4/k')+o(1)$ and
$K(k')=\pi/2+O(k'^2)$ then give
$s_1(t)=\tfrac13\log(4/t)+o(1)$, which proves
\eqref{eq:s-chainlog} with $c_0=(\log4)/3$.

\paragraph{A shifted-measure proof of the grouped lattice limit.}
Set $\epsilon=a^2$ and push the tangential
measure forward by $u=\sin^2\theta_2+\sin^2\theta_3$. Each
single sine square has density
$[\pi\sqrt{u(1-u)}]^{-1}$ on $(0,1)$. Its convolution satisfies
\begin{equation}
 w(u)=\frac1{\pi^2}\int_0^u
 \frac{\dd v}{\sqrt{v(u-v)(1-v)(1-u+v)}}
 =\frac1\pi+O(u),\qquad 0<u<1.
 \label{eq:s-density}
\end{equation}
It is smooth near zero; its integrable singularity at $u=1$
is separated from the critical region. Define
$H(u)=s_1(\sqrt u)$ and
$\Delta_\epsilon H(u)=\sum_s\eta_sH(u+\epsilon x_s)$.
The exact expression is
\[
 \frac{S_{\rm gr}}{\mathcal A}
 =\frac2\epsilon\int_0^2w(u)\Delta_\epsilon H(u)\dd u.
\]
Choose a smooth cutoff supported near zero. On its complement
the entropy is gapped and smooth, so the grouped integral is
$O(\epsilon^2)$ by the two mass moments.

For the near-zero remainder $r(u)=H(u)+(1/6)\log u$, let $v$
be the cutoff times $w$ and subtract $r(0)=c_0$. Shifting
$y=u+\epsilon x_s$ gives
\[
 \int_0^\infty[r(y)-r(0)]
 \sum_s\eta_s v(y-\epsilon x_s)
                 \id_{y\ge\epsilon x_s}\dd y.
\]
Above $\epsilon\max_sx_s$, Taylor expansion of the smooth
weight, rather than of the entropy, gives $O(\epsilon^2)$ after
the two moment cancellations. In the remaining interval, whose
length is $O(\epsilon)$, continuity of $r$ gives $o(\epsilon)$.
The remainder divided by $\epsilon$ therefore vanishes.

For the logarithm with constant weight,
\[
 \frac1\epsilon\int_0^\delta
       \Delta_\epsilon\log u\dd u\longrightarrow-J.
\]
Evaluate $(u+\epsilon x)\log(u+\epsilon x)-(u+\epsilon x)$
at the endpoints and apply the two mass moments. The $O(u)$
part of $w$ contributes $O(\epsilon^2|\log\epsilon|)$ before
division by $\epsilon$: the rectangle difference of the logarithm
is bounded by a constant times
$\epsilon^2/(u+\epsilon x_0)^2$. These steps give
\begin{equation}
 \lim_{a\to0}\frac{S_{\rm gr}(a)}{\mathcal A}
 =\frac1{3\pi}\int_0^{h_1}\!\dd u\int_0^{h_2}\!\dd v\,
       \frac1{x_0+u+v}
 =\frac{J}{3\pi}.
 \label{eq:s-group-limit}
\end{equation}
Equivalently one packet has
$F_a(x)=c_{-2}/a^2+(x/(3\pi))\log(a^2x)+c_2x+o(1)$
with a cut-regulator-dependent $c_{-2},c_2$. Both the area divergence
and the $x\log a$ term cancel in the grouped expression; the
linear analytic mass term cancels as well. No finite area counterterm
is fitted.

\section*{4. Direct continuum cone normalization}
For one four-component Dirac valley on
$C_n\times\mathbb R^2$, the same gluing calculation uses
\[
 \tr_4(R/6-R/4)=-R/3,
 \qquad\int R=4\pi(1-n)\mathcal A.
\]
The surface heat trace is
$(n-1)\mathcal A/(12\pi t)$ to first order in $n-1$.
The Grassmann determinant therefore gives the cut area density
\begin{equation}
 \frac{S_D(m)}{\mathcal A}
 =\frac1{24\pi}\int_{\epsilon^2}^{\infty}
          \frac{\dd t}{t^2}e^{-m^2t}.
 \label{eq:s-cone}
\end{equation}
This determines the continuum mass-dependent area term; its
power-divergent coefficient is not identified with the sine cutoff's
coefficient. Grouping first removes the lower cutoff:
\begin{align}
 \sum_s\eta_s\frac{S_D(m_s)}{\mathcal A}
 &=\frac1{24\pi}\int_0^{\infty}\frac{\dd t}{t^2}
 e^{-x_0t}(1-e^{-h_1t})(1-e^{-h_2t})\notag\\
 &=\frac{J}{24\pi}.
 \label{eq:s-cone-group}
\end{align}
The second line follows by writing each exponential difference as
an integral over its positive mass increment and using Tonelli.
Eight valleys multiply it to \eqref{eq:s-group-limit}. The
half-chain Toeplitz/CTM spectrum and shifted lattice measure give
that result without a cone input. The cone computation is therefore
a distinct check of the coefficient, rather than a change of
dimension in the same proof.

\section*{5. Comparison with the actual vacuum stiffness}
The exact finite-spacing, thermodynamic small-momentum coefficient
of the chosen polar law is
\[
 k_a(1)=\frac3{a^2}\int_{[-\pi,\pi]^3}
 \frac{\dd^3\theta}{(2\pi)^3}\cos^2\theta_1
 \sum_s\frac{\eta_s}{\sqrt{\sum_i\sin^2\theta_i+a^2x_s}},
 \qquad k_a(1)\longrightarrow\frac{3J}{\pi^2}.
\]
$G$ is defined by the Cavendish measurement. The single relation
$k_a(1)=1/(2\pi G)$ is a condition on the inventory at the chosen
spacing. The proposed area density is consequently
$(\pi/2)k_a(1)$, whose continuum value is $3J/(2\pi)$.
The graded native cut coefficient in \eqref{eq:s-group-limit} is
smaller by the factor $2/9$. The common moment cancellations alone
do not make the two finite parts equal.

The cut is fixed by the flat negative-band projector. The vacuum
coefficient is the response of the \emph{chosen inhomogeneous polar
allocation}, including its nonminimal gradient placement and
bosonic frequency allocation. Equality with a covariantly coupled
Dirac determinant's induced Newton term cannot be imported merely
from the shared mass rectangle. The exact finite-$a$ cut remains
\eqref{eq:s-exact}; no all-spacing identity or error enclosure is
claimed for it.

\paragraph{Primary continuum coefficient checks.}
These checks use no entropy result to infer the polar response.
For a minimally coupled real scalar, Demers--Lafrance--Myers
obtain the regulated area term $B\mathcal A/(48\pi)$ and the
same $B/(12\pi)$ contribution to $G^{-1}$, using a common
covariant regulator. Their nonminimal-coupling example changes the
Newton coefficient while leaving the stated mode-counting term
unchanged. Their result does not identify an arbitrary vacuum
response with a cut entropy.

Frolov--Fursaev--Zelnikov give the Dirac contribution
$G_{\rm ind}^{-1}=\sum_d m_d^2\log m_d^2/(6\pi)$ and the
regulated one-Dirac area coefficient
$(24\pi)^{-1}\int\dd t\,t^{-2}e^{-m_d^2t}$.
The latter agrees with \eqref{eq:s-cone}; applying our signed
rectangle gives $J/(24\pi)$ per valley. Their induced-gravity
model also has curvature-coupled scalars and a horizon contact term:
the total black-hole entropy is the statistical entropy minus
$2\pi\sum_s\xi_s\int_\Sigma\langle\phi_s^2\rangle$.
That subtraction and their field constraints are not ingredients
of the native regional CAR cut defined here. Applying the same
signed rectangle to eight covariant Dirac valleys per packet gives
the comparison coefficient
$G_{\rm ind}^{-1}=4J/(3\pi)$, whereas the declared polar match
gives $G^{-1}\to6J/\pi$. This is the same factor $2/9$, not a
second calibration of $H_*$.

\paragraph{Limits of this conclusion.}
The result is a planar area density of the flat prescribed sea.
Smooth curved cuts can carry additional curvature terms and finite
shape dependence; the argument is not a full entropy formula for
a finite sphere or a physical apparent horizon. The lower corner
is strictly positive. A massless-corner coefficient limit can be
taken separately; it does not give a gapped polar inverse.
Positivity of the limit does not make the signed combination the
ordinary entropy of that regional algebra or establish a thermal
first law.

For a homogeneous positive ruler $\psi=P$ and positive constant
lapse, the same calculation replaces $a$ by $a_{\rm prop}=aP^2$
and $\mathcal A$ by $\mathcal A_{\rm prop}=P^4\mathcal A_{\rm coord}$.
Thus the limiting coefficient per proper area remains $J/(3\pi)$.
Changing only the area denominator while keeping the unit-ruler
projector would be a different calculation. This homogeneous
rescaling does not establish a curved-cut or trapping-surface
entropy law.

\section*{SOURCES and independence record}
\begin{enumerate}[leftmargin=*]
\item Governing task and history:
\path{<internal>/GOAL_alt.md};
\path{<internal>/JUDGMENT_web40c.md},
sections 12--13;
\path{<internal>/NOTES_shift_and_spacing_window.md};
\path{<internal>/GOAL_lane_z.md}.
The signed cut/equality is a question to test. The minus sign attached
to $\sum\eta_sx_s\log x_s$ in one sentence of judgment 13(1) is
inconsistent with its positive rectangle integral;
\eqref{eq:s-J} supplies the explicit correction.
\item The fixed law and protected sea:
\path{<internal>/model_and_hypotheses.tex};
\path{<internal>/strand_A_lattice.tex};
\path{<internal>/review_A.tex},
polar block diagonalization, site inventory and protected vacuum.
These are antecedent model sources, not either lane's Strand S.
\item Actual vacuum normalization:
\path{<internal>/rederivation_final3_vacuum_surface.tex};
\path{<internal>/redteam_final3_vacuum_surface.tex};
\path{<internal>/p19_cross_lane_agreement.tex}.
The full Brillouin zone contains the eight valleys once.
\item Exact one-interface chain input: I.~Peschel,
\emph{On the entanglement entropy for a XY spin chain},
J.~Stat.~Mech. (2004) P12005,
\href{https://arxiv.org/abs/cond-mat/0410416}{arXiv:cond-mat/0410416},
equations (1), (6)--(8), (13) and the one-interface/two-interface
paragraph following (14), pages 1--3.
Actual primary text read:
\path{<internal>/Peschel_0410416.txt}.
Its printed range $0\le\gamma\le1$ is respected by the Hardy-space
map, rather than extrapolated to a pure-pairing parameter limit.
A.~R.~Its, B.-Q.~Jin and V.~E.~Korepin,
\emph{Entropy of XY Spin Chain and Block Toeplitz Determinants},
\href{https://arxiv.org/abs/quant-ph/0606178}{arXiv:quant-ph/0606178v3},
equations (23)--(31), pages 7--9, in
\path{<internal>/IJK_0606178.txt},
was checked for boundary normalization and scope. Its equation (28)
uses an ordered regularized double limit; the following paragraph
calls interchange conjectural. The proof here uses Peschel's
half-chain CTM spectrum, not that interchange.
\item Continuum checks only: J.-G.~Demers, R.~Lafrance and
R.~C.~Myers,
\emph{Black hole entropy without brick walls},
\href{https://arxiv.org/abs/gr-qc/9503003}{arXiv:gr-qc/9503003},
equations (6)--(10), (18)--(23), (28) and section 4.4;
\emph{Black Hole Entropy and Renormalization},
\href{https://arxiv.org/abs/gr-qc/9507042}{arXiv:gr-qc/9507042},
equations (2.3)--(2.8), (3.5)--(3.9).
Actual texts:
\path{<internal>/DLM_9503003.txt} and
\path{<internal>/DLM_9507042.txt}.
They concern covariantly regulated scalar fields and do not supply
the native sine cut or its polar stiffness.
V.~P.~Frolov, D.~V.~Fursaev and A.~I.~Zelnikov,
\emph{Statistical Origin of Black Hole Entropy in Induced Gravity},
\href{https://arxiv.org/abs/hep-th/9607104}{arXiv:hep-th/9607104},
equations (2.1)--(2.19), (3.2)--(3.12), (4.26)--(4.30), in
\path{<internal>/FFZ_9607104.txt}.
The Dirac factor, common regulator and scalar contact subtraction
are retained at their stated scopes.
\item Independence and actual work: the substantive exact
dimensional reduction, cone coefficient, rectangle grouping and
stiffness comparison were written before reading these continuum
coefficient checks or any other lane's Strand S work. No icloud
Strand S derivation, task result or proposed route, and no other
lane's numerical Strand S result, has been read. The independent
RED proposed the Hardy-space map; the author reconstructed its
unitarity, branch-point map, parameter inequalities and spectrum
factor in the actual proof above. The shifted-measure limit was
also reconstructed directly. The agreeing mathematical routes are
the exact Toeplitz/CTM reduction with the lattice measure and the
grouped Dirac-cone calculation. No scientific code, postprocessing
or numerical entropy evaluation has been run by the author.
\end{enumerate}
% The exact Hardy-space map keeps the CTM parameters inside the primary's
% stated domain. The shifted-measure proof does not differentiate the
% critical entropy remainder. No numerical specification or run is needed.
\end{document}
